%% ──────────────────────────────────────────────
%%  فصل ۱۵ – بالاتر از سیاهی رنگی نیست:
%%  نجات‌طلبی فداکارانه و مسئولیت نسلی
%%  فایل: chapters/ch15-sacrificial-liberation.tex
%% ──────────────────────────────────────────────
\chapter{بالاتر از سیاهی رنگی نیست: نجات‌طلبی فداکارانه و مسئولیت نسلی}
\label{ch:sacrificial-liberation}

%% ── چکیده‌ی فصل ──
\begin{tcolorbox}[mybox, title={چکیده‌ی فصل}]
فصل‌های پیشین این کتاب استدلال کردند که نجات‌طلبی — امیدبستن به نیروی خارجی برای رهایی — در اکثر موارد تاریخی شکست خورده و نتایجی از بی‌ثباتی تا فاجعه به‌بار آورده است. اما آیا این تحلیل، تمام داستان است؟ این فصل عمداً نقش \concept{مدافع شیطان} (\lat{Devil's Advocate}) را ایفا می‌کند و قوی‌ترین استدلال‌های ممکن را \textbf{به نفع} نجات‌طلبی — آنچه \concept{نجات‌طلبی فداکارانه} می‌نامیم — ارائه می‌دهد. پرسش محوری این است: \textbf{وقتی انسداد سیاسی مطلق است و هیچ راه داخلی‌ای وجود ندارد، آیا پذیرش هزینه‌ی مداخله‌ی خارجی — حتی به قیمت یک یا دو نسل — برای رهایی نسل‌های آینده موجه نیست؟} اگر «بالاتر از سیاهی رنگی نیست»، آیا هر تغییری — حتی پرهزینه — بهتر از تداوم وضع موجود نیست؟ این فصل ابتدا چارچوب‌های مختلف پاسخ‌گویی به این پرسش را معرفی می‌کند، سپس شواهد تاریخی و نظری هر رویکرد را بررسی و در نهایت ارزیابی می‌نماید.
\end{tcolorbox}

%% ================================================================
\section{درآمد: پرسشی که نباید نادیده گرفت}
\label{sec:ch15-intro}
%% ================================================================

\needspace{6\baselineskip}
نویسنده‌ی این کتاب در فصل‌های ۱ تا ۱۴ استدلال کرد که نجات‌طلبی معمولاً آرزواندیشی‌ای خطرناک است. اما صداقت فکری ایجاب می‌کند که قوی‌ترین استدلال مخالف نیز شنیده شود — استدلالی که از دل تاریک‌ترین تجربه‌های بشری برمی‌آید:

\begin{tcolorbox}[quotebox, title={صدای دیگر}]
«شما از نروژ یا سوئیس درباره‌ی عاملیت جمعی و مقاومت مدنی سخن می‌گویید. اما من از آشویتس سخن می‌گویم. از حلبچه. از سربرنیتسا. از رواندا. در آشویتس هیچ مقاومت مدنی‌ای کارساز نبود. هیچ نهاد مدنی مستقلی نمی‌توانست کوره‌های آدم‌سوزی را متوقف کند. تنها ارتش سرخ و متفقین توانستند آشویتس را آزاد کنند. آیا یهودیانی که از متفقین تقاضای بمباران خطوط ریلی به آشویتس را داشتند، "نجات‌طلب" بودند؟ آیا آرزواندیشی می‌کردند؟ یا واقع‌بینانه‌ترین ارزیابی ممکن را از وضعیت‌شان داشتند؟»\\
— مضمونی بازسازی‌شده از استدلال‌های منتقدان ضدمداخله‌گرایی مطلق
\end{tcolorbox}

این صدا را نمی‌توان با آمار «۸۴٪ مداخلات شکست خورده‌اند» ساکت کرد. زیرا برای کسی که در تاریکی مطلق زندگی می‌کند، حتی شمعی لرزان بهتر از هیچ است. پرسش واقعی این نیست که «آیا مداخله خطر دارد؟» — البته که دارد. پرسش واقعی این است: \textbf{«آیا هزینه‌ی مداخله، کمتر از هزینه‌ی عدم مداخله است؟»}

\begin{tcolorbox}[defbox, title={مفهوم کلیدی: نجات‌طلبی فداکارانه}]
\concept{نجات‌طلبی فداکارانه} (\lat{Sacrificial Liberation-Seeking}) رویکردی است که در آن یک ملت (یا بخشی از آن) آگاهانه و با چشم باز هزینه‌های مداخله‌ی خارجی — شامل خشونت، بی‌ثباتی، وابستگی موقت، و حتی تلفات نسلی — را می‌پذیرد، نه از سر آرزواندیشی یا خام‌اندیشی، بلکه بر اساس \concept{محاسبه‌ی عقلانی} مبنی بر اینکه:
\begin{enumerate}
\item وضع موجود (انسداد مطلق) بدتر از بدترین پیامد مداخله است
\item راه‌های داخلی تغییر واقعاً و نه فرضاً مسدود شده‌اند
\item هزینه‌ی پرداخت‌شده توسط نسل فعلی، سرمایه‌گذاری برای آزادی نسل‌های آینده است
\item جایگزین (عدم اقدام)، نه «صلح» بلکه \concept{خشونت ساختاری مداوم} است
\end{enumerate}
تفاوت بنیادین نجات‌طلبی فداکارانه با نجات‌طلبی ساده (بررسی‌شده در فصل‌های قبل): \textbf{آگاهی از هزینه‌ها + پذیرش آگاهانه + افق نسلی}.
\end{tcolorbox}


%% ================================================================
\section{پنج چارچوب پاسخ‌گویی: از کجا به مسئله نگاه کنیم؟}
\label{sec:ch15-frameworks}
%% ================================================================

\needspace{6\baselineskip}
پاسخ به پرسش «آیا نجات‌طلبی فداکارانه موجه است؟» عمیقاً به چارچوب تحلیلی‌ای بستگی دارد که از آن به مسئله نگاه می‌کنیم. پنج چارچوب اصلی وجود دارد:

\subsection{چارچوب اول: فایده‌گرایی نسلی}
\label{subsec:ch15-utilitarianism}

\concept{فایده‌گرایی نسلی} (\lat{Intergenerational Utilitarianism}) — مبتنی بر سنت فکری \thinker{جرمی بنتام} (\lat{Jeremy Bentham}) و \thinker{جان استوارت میل} (\lat{John Stuart Mill}) — اقدام درست را اقدامی می‌داند که بیشترین خیر (یا کمترین رنج) را برای بیشترین تعداد افراد، \textbf{در طول زمان}، تولید کند.\footnote{\lr{Parfit, Derek. \textit{Reasons and Persons}. Oxford: Clarendon Press, 1984, pp. 351–379.}}

\begin{tcolorbox}[wavebox, title={استدلال فایده‌گرایانه به نفع نجات‌طلبی فداکارانه}]
فرض کنید:
\begin{itemize}
\item یک دیکتاتوری مطلق سالانه ۱۰٬۰۰۰ نفر را می‌کشد و ده‌ها میلیون را در فقر و سرکوب نگه می‌دارد
\item هیچ راه داخلی‌ای برای تغییر وجود ندارد (سرکوب کامل)
\item مداخله‌ی خارجی ممکن است ۱۰۰٬۰۰۰ نفر تلفات فوری داشته باشد و ۱۰ سال بی‌ثباتی ایجاد کند
\item اما پس از آن، نظامی دموکراتیک (یا حداقل کمتر سرکوبگر) برقرار شود
\end{itemize}

\textbf{محاسبه:}
\begin{itemize}
\item هزینه‌ی تداوم وضع موجود (۵۰ سال): ۵۰۰٬۰۰۰ کشته + ده‌ها میلیون سرکوب‌شده
\item هزینه‌ی مداخله: ۱۰۰٬۰۰۰ تلفات + ۱۰ سال بی‌ثباتی
\item سود مداخله (اگر موفق شود): آزادی نسل‌های آینده
\end{itemize}

اگر \textbf{احتمال موفقیت} حداقل ۲۰–۳۰٪ باشد، فایده‌ی انتظاری (\lat{Expected Utility}) مداخله ممکن است از عدم مداخله بیشتر باشد.
\end{tcolorbox}

\thinker{درک پارفیت} (\lat{Derek Parfit}) در کتاب \textit{دلایل و اشخاص} استدلال کرده است که ما مسئولیت اخلاقی نسبت به نسل‌های آینده داریم — و این مسئولیت ممکن است ایجاب کند که نسل فعلی هزینه‌هایی را بپذیرد که به نفع آیندگان است.\footnote{\lr{Parfit, Derek. \textit{Reasons and Persons}, op. cit., pp. 351–379.}}

\subsection{چارچوب دوم: وظیفه‌گرایی کانتی}
\label{subsec:ch15-kantian}

\concept{وظیفه‌گرایی کانتی} (\lat{Kantian Deontology}) — مبتنی بر فلسفه‌ی \thinker{ایمانوئل کانت} (\lat{Immanuel Kant}) — بر حقوق بنیادین انسان تأکید دارد: هیچ‌کس نباید صرفاً ابزار اهداف دیگران باشد (\concept{امر مطلق} / \lat{Categorical Imperative}).

\begin{tcolorbox}[wavebox, title={استدلال کانتی: دو تفسیر متضاد}]
\textbf{تفسیر ضد مداخله (غالب در این کتاب):}
استفاده از ملت‌ها به‌عنوان «ابزار» اهداف ژئوپولیتیک مداخله‌گر، نقض امر مطلق کانتی است. مردم عراق «وسیله»ی اثبات قدرت آمریکا شدند — نه «غایت» فی‌نفسه.

\textbf{تفسیر حامی مداخله (بررسی‌شده در این فصل):}
تداوم دیکتاتوری نیز نقض امر مطلق است: دیکتاتور مردم خود را «ابزار» حفظ قدرتش می‌کند. اگر مداخله‌ی خارجی تنها راه بازگرداندن \concept{کرامت انسانی} (\lat{Human Dignity}) باشد، \textit{عدم} مداخله خود نقض وظیفه‌ی اخلاقی است.

\thinker{آلن بیوکنن} (\lat{Allen Buchanan})، فیلسوف سیاسی، استدلال کرده است که «حق حاکمیت» مشروط به رعایت حقوق اساسی شهروندان است — و دولتی که این حقوق را نقض کند، «حق» حاکمیت را از دست می‌دهد.\footnote{\lr{Buchanan, Allen. \textit{Justice, Legitimacy, and Self-Determination: Moral Foundations for International Law}. Oxford: Oxford University Press, 2004, pp. 233–270.}}
\end{tcolorbox}

\subsection{چارچوب سوم: عدالت بین‌نسلی}
\label{subsec:ch15-intergenerational}

\concept{عدالت بین‌نسلی} (\lat{Intergenerational Justice}) پرسش می‌کند: آیا نسل فعلی حق دارد — یا حتی \textit{وظیفه} دارد — برای آزادی نسل‌های آینده هزینه بپردازد؟

\thinker{جان رالز} (\lat{John Rawls}) در \textit{نظریه‌ی عدالت} مفهوم \concept{اصل پس‌انداز عادلانه} (\lat{Just Savings Principle}) را مطرح کرده:

---

## Conversation 2
*Model: claude-opus-4-6*

### Request

این فصل را کامل کن

‍‍